%! Piecewise \& Step Functions — Unit 2, Lesson 2.4 — Student Packet Cover Sheet
\documentclass[10pt]{article}
\usepackage{algebra2-article}
\usepackage{algebra2-boxes}
\usepackage{ltablex}
\keepXColumns

\begin{document}

% Full-bleed forest banner
\begin{tikzpicture}[remember picture, overlay]
  \fill[forest] (current page.north west)
    rectangle ([yshift=-0.9in]current page.north east);
\end{tikzpicture}
\vspace{-0.8in}

\begin{center}
  {\color{white}\textbf{\LARGE Algebra 2: Shepherd}}\\[4pt]
  {\color{forestmid}\large\textbf{Unit 2 \quad Linear Functions}}\\[6pt]
  {\color{white}\textbf{\large Lesson 2.4 \quad Piecewise \& Step Functions}}
\end{center}

\vspace{0.22in}
\namedateperiod
\vspace{0.18in}

\begin{learningtargetbox}
\small
\begin{enumerate}[leftmargin=1.5em, itemsep=5pt]
  \item \ldots{} read \textbf{piecewise notation} and \textbf{evaluate} $f(x)$ by choosing the piece whose condition $x$ satisfies (minding the boundary).
  \item \ldots{} read a piecewise graph, using \textbf{open} and \textbf{closed} endpoints to find where a piece starts and stops and to spot a \textbf{discontinuity} (jump).
  \item \ldots{} write the \textbf{absolute value} function $|x|$ (and $|x-h|$) as a \textbf{two-piece} linear rule.
  \item \ldots{} evaluate and read the \textbf{greatest-integer (step) function} $\lfloor x\rfloor$, and read domain, range, and increasing/decreasing/\textbf{constant} intervals.
\end{enumerate}
\end{learningtargetbox}

\vspace{0.14in}

\begin{tocbox}
\small
\renewcommand{\arraystretch}{1.5}
\begin{tabularx}{\linewidth}{@{} c l X r @{}}
\rowcolor{forestbg}
\textbf{\#} & \textbf{Component} & \textbf{Description} & \textbf{Score} \\
\midrule
1 & Warm-Up        & Spiral review: evaluate a rule, $|x|$ as two lines, open/closed endpoints & \blank{1.2cm} \\
2 & Guided Notes   & Piecewise notation, evaluating, reading graphs, $|x|$ as pieces, step functions & \blank{1.2cm} \\
3 & Group Activity & \emph{Pieces, Graphs \& Staircases} --- evaluate, read graphs, model, by tier & \blank{1.2cm} \\
4 & Exit Ticket    & Quick check: evaluate across a boundary, read a graph, greatest integer & \blank{1.2cm} \\
5 & Homework       & Mixed practice: evaluate, match, read, write as pieces, step function, model & \blank{1.2cm} \\
\midrule
  & \multicolumn{2}{r}{\textbf{Total}} & \blank{1.2cm} \\
\end{tabularx}
\end{tocbox}

\vspace{0.10in}

\begin{remindbox}
\small A \textbf{piecewise-defined function} follows different rules on different parts of its domain,
written with a brace, e.g.\ $f(x)=\{\,2x+1 \text{ if } x<2;\ 5-x \text{ if } x\ge 2\,\}$. To
\textbf{evaluate}, first decide \emph{which piece} the input belongs to (the boundary's $<$ vs.\ $\le$
tells you which piece owns it), then apply that rule. On a graph, a \textbf{closed} dot $\bullet$
includes its endpoint and an \textbf{open} dot $\circ$ excludes it; when the two pieces disagree at a
boundary the graph \emph{jumps} (a \textbf{discontinuity}). Absolute value is piecewise
($|x|=\{\,x \text{ if } x\ge 0;\ -x \text{ if } x<0\,\}$), and the \textbf{greatest-integer function}
$\lfloor x\rfloor$ (round \emph{down}) is a \textbf{step} function: constant on each interval, then a
jump.
\end{remindbox}

\end{document}
